\scanpage{175}
\mspage{10}
% archive sheet number: 174
2.2. \uline{Réduction de} 1.3 \uline{au cas} $X$ \deleted{géom.} \uline{affine lisse de dim 1}.

Soit $U_K$ un ouvert affine \deleted{\unsure{connexe|non vide}} de $X_K$, \deleted{Alors} contenant le pt \deleted{base géom.} $x$. Alors $\pi_1'(U_{\bar K},\bar x)\mathrel{\text{---}}\pi_1(X_{\bar K},\bar x)$ est surjectif et compatible avec les op. de $I$, donc de même pour $\pi_1^{\unsure{(p')|(p)}}(U_{\bar K},\bar x)\mathrel{\text{---}}\pi_1^{(p')}(X_{\bar K},\bar x)$, donc on est ramené au cas de $X$ affine, a fortiori quasi-projectif. La réduction \inserted{au cas} où de plus $\dim X=1$ se fait alors à coups de sections hyperplanes génériques, comme en II 3. \deleted{\unsure{Comme}}

2.3. \uline{Réduction de} 1.3 \uline{au cas de la situation 1.4, avec} \deleted{\illegible{1}} $\uline{\dim} X_K=1$.
\inserted{En vertu de I, on peut supposer \unsure{maint.|normal} $S$ complet.}
Considérons le modèle complet $\bar X_K$ de $X_K$. Quitte à faire une extension radicielle de $K$, on peut supposer $\bar X_K$ géom. normal, i.e. lisse.
De plus, un théorème non trivial, qui sera vu dans un exposé ultérieur (\quad), nous apprend que, quitte à \deleted{\unsure{remplacer}} passer à une extension finie de $K$, il existe un \deleted{\unsure{mod}} schéma projectif $\bar X$ sur $S$, ayant $\bar X_K$ comme fibre générique, et tel que $\bar X_s$ soit \deleted{géométriquement réduit} un diviseur à croisements normaux dans $\bar X$, $X_s$ étant de plus réduit et son normalisé étant \deleted{géom.} \uline{lisse} sur $\unsure{K|k}$.

\margin{corps résiduel parfait ? Mais on s'y ramène pour prouver 1.3, \deleted{\unsure{déjà fait}} \inserted{par réduction} à un ss groupe \deleted{\unsure{distingué}} \inserted{ouvert} de $I$, cf I\unsure{...}}

Agrandissant encore $K$, on peut supposer que les pts de $\bar X_K-X_K$ sont rationnels sur $K$, donc \inserted{($\bar X$ étant propre sur $S$)} sont définis par des sections $g_i$ de $\bar X$ sur $S$, $1\leq i\leq r$. La réduction est
